\paragraph{Cluster A Step 14: real anomaly enrichment.}

This step replaces the toy additive anomaly score with exact rational anomaly
coefficients on a finite curated set of real chiral multiplet candidates.  The
reference one-generation content is
\[
Q:(3,2)_{1/6},\quad u^c:(\bar 3,1)_{-2/3},\quad
d^c:(\bar 3,1)_{1/3},\quad L:(1,2)_{-1/2},\quad e^c:(1,1)_1 .
\]

The computed coefficients for this content are all zero:
\[
[SU(3)]^3=0,\quad [SU(3)]^2U(1)=0,\quad [SU(2)]^2U(1)=0,\quad
U(1)^3=0,\quad U(1)\text{-grav}=0 ,
\]
and the number of \(SU(2)\) doublets is even, so the Witten global anomaly is
absent. This is known physics limit recovery, used here as a faithfulness
calibration.

For fixed one-generation representations, the anomaly equations are
\[
2Y_Q+Y_{u^c}+Y_{d^c}=0,\qquad 3Y_Q+Y_L=0,
\]
\[
6Y_Q+3Y_{u^c}+3Y_{d^c}+2Y_L+Y_{e^c}=0,
\]
\[
6Y_Q^3+3Y_{u^c}^3+3Y_{d^c}^3+2Y_L^3+Y_{e^c}^3=0 .
\]
Solving gives
\[
Y_L=-3Y_Q,\qquad Y_{e^c}=6Y_Q,\qquad
\{Y_{u^c},Y_{d^c}\}=\{-4Y_Q,2Y_Q\}.
\]
With the usual up/down labeling and \(Y_Q=1/6\), this recovers the familiar
hypercharge pattern and \(Q_{\mathrm{proton}}=-Q_e\).

The finite candidate space contains anomaly-free non-reference survivors, so
real anomaly freedom remains necessary but not sufficient.  The scoped next
consequence is a minimal irreducible anomaly-cancellation support test over a
larger real candidate enumeration.

\paragraph{Grade.}
Finite-carrier diagnostic and known-physics limit recovery.  No physical SM
value or mechanism is supplied, and frame-transfer remains open.
